\section{Development priorities and versioned guarantees}\label{sec:evolution}

Four identifiers must remain explicit: manuscript version, software version and source identity, certificate schema, and semantic contract. Manuscript 2.0 accompanies the preparation of software \texttt{0.2.0a1}; the coordinatewise schema and contract remain unchanged, while research capsules preserve their individual schemas and source fingerprints. Lean additionally pins its compiler version. A proposed release tag is not evidence that a remote publication has occurred.

The release manifest maps K1--K10 to code, saved proof data, and validation. The manuscript has no invented DOI. If it is deposited later, a version DOI should identify this exact text and a concept DOI the evolving record. The public repository should preserve the earlier paper/software map and identify the newly released commit after publication.

The next research priority is the complete joint-carrier argument for Graal \texttt{create}. It requires deriving helper preconditions from reachable caller states, preserving temporary masks and empty outcomes, and proving the stated stabilization behavior. In particular, the lower sign condition must be established where it is used; it cannot be assumed because it made a standalone helper theorem true.

The next formalization step can reuse the general gluing theorem for the 49-state descending certificate, then address the outer order cases and source correspondence. Extending the algebraic Lean carrier beyond $\mathcal P(\{0,1\})$ is useful when a concrete next proof requires it. Formalizing the whole series at once is not a prerequisite for checking a stable applied fragment.

The global algorithmic challenge is observation inference. Current producers automatically complete rows, search closed carriers, and check proposed certificates within designed observation families. They do not yet derive a suitable finite family for any source program. A stronger engine should derive candidate observations and their sufficient domains from consumer obligations, retain separating countermodels when a factor loses essential information, and report exhaustion explicitly.

Wider application also requires target adapters matching actual language preconditions, integration of the research interfaces into a stable public API, and controlled comparisons on new external sources. Coverage, certificate size, discovery effort, replay effort, native correspondence, and proof-assistant coverage are distinct progress measures. Each should retain its own denominator and failure outcomes.
